%ID BT_6_1_1_2_01
Dùng quy ước $f=a_0+\sum_{n\ge1}(a_n\cos k_nx+b_n\sin k_nx)$, $k_n=n\pi/L$, $a_0=(2L)^{-1}\int_{-L}^Lf$, $a_n=L^{-1}\int f\cos$, $b_n=L^{-1}\int f\sin$.
(i) $a_0=\sin(\alpha L)/(\alpha L)$, $a_n=2(-1)^n\alpha\sin(\alpha L)/[L(\alpha^2-k_n^2)]$, $b_n=0$. Ở $\alpha=0$ hoặc cộng hưởng, lấy giới hạn (hàm hằng hoặc một mode cos). Chuỗi hội tụ đều đến phép kéo dài liên tục chu kỳ $2L$.
(ii) $b_n=2(1-(-1)^n)/(n\pi)$, các $a_n=0$. Chuỗi hội tụ đến sign trong miền trừ điểm 0, về 0 tại 0 và các điểm nối $\pm L$; không hội tụ đều qua bước nhảy.
(iii) $a_0=\sinh(\alpha L)/(\alpha L)$,
\[a_n=\frac{2\alpha(-1)^n\sinh(\alpha L)}{L(\alpha^2+k_n^2)},\quad b_n=\frac{2k_n(-1)^{n+1}\sinh(\alpha L)}{L(\alpha^2+k_n^2)}.\]
Với $\alpha=0$, $f=1$. Với $\alpha\ne0$, chuỗi về $e^{\alpha x}$ trong $(-L,L)$, và $\cosh\alpha L$ tại mép; hội tụ đều trên compact nội, không trên toàn chu kỳ.
(iv) Phần ảo của $\log(1+re^{ix})$ cho $f=\sum_{n\ge1}(-1)^{n+1}r^n\sin nx/n$.
(v) Hai lần phần thực cho $f=2\sum_{n\ge1}(-1)^{n+1}r^n\cos nx/n$.
(vi) Nhân Poisson cho $f=(1-r^2)^{-1}[1+2\sum_{n\ge1}r^n\cos nx]$. Ở (iv)--(vi), $0<r<1$, chuỗi hội tụ tuyệt đối và đều theo tiêu chuẩn Weierstrass.
%ID BT_6_1_1_2_02
(i) Khai triển Fourier $\cos(\alpha x)$ trên $[-\pi,\pi]$ và thế $x=\pi$ cho
\[\sum_{n\in\mathbb Z}\frac1{\alpha^2-n^2}=\frac\pi\alpha\cot\pi\alpha,\quad0<|\alpha|<1.\]
Điều kiện $|\alpha|<1$ trong đề còn phải loại $\alpha=0$ vì số hạng $n=0$ không xác định.
(ii) Khai triển $\cosh(\alpha x)$ (tích phân thực tương tự) và thế $x=\pi$ cho
\[\sum_{n\in\mathbb Z}\frac1{\alpha^2+n^2}=\frac\pi\alpha\coth\pi\alpha,\quad\alpha\ne0.\]
Các chuỗi hội tụ tuyệt đối vì số hạng $O(n^{-2})$, nên phép gom hai phía hợp lệ.
%ID BT_6_1_1_2_03
(i) $\sum\sin nx/n=(\pi-x)/2$ với $0<x<2\pi$, kéo dài chu kỳ $2\pi$; bằng 0 tại $x\in2\pi\mathbb Z$. Suy từ Fourier của hàm răng cưa hoặc phần ảo $-\log(1-re^{ix})$ rồi $r\uparrow1$.
(ii) Tích phân chuỗi (i), ấn định hằng bằng trung bình 0, cho $\sum\cos nx/n^2=\pi^2/6-\pi x/2+x^2/4$ trên $[0,2\pi]$, kéo dài tuần hoàn. Chuỗi hội tụ đều và tuyệt đối.
(iii) Tổng chuỗi hình học với $z=-e^{ix}/2$:
\[\sum_{n\ge1}(-1)^n2^{-n}\sin nx=\operatorname{Im}\frac z{1-z}=-\frac{2\sin x}{5+4\cos x}.\]
Chuỗi hội tụ tuyệt đối và đều.
%ID BT_6_1_1_2_04
(i) Tích phân từng phần $k$ lần, các hạng biên mất vì tuần hoàn; hệ số phức $\widehat f(n)=(in)^{-k}\widehat{f^{(k)}}(n)$. Bổ đề Riemann--Lebesgue cho $\widehat{f^{(k)}}(n)\to0$, nên $n^k(|a_n|+|b_n|)\to0$.
(ii) Giả thiết cho các chuỗi đạo hàm đến cấp $k$ được chặn bởi $M\sum n^{-2}$. Vì vậy tổng Fourier có đại diện $g\in C^k$, và $f=g$ hầu khắp nơi do tính duy nhất hệ số Fourier. Đề chỉ giả thiết $f$ khả tích nên không thể kết luận $f\in C^k$ tại mọi điểm: sửa một giá trị của $g$ không đổi bất kỳ hệ số nào nhưng có thể phá liên tục. Kết luận nguyên văn cần thêm giả thiết $f$ liên tục, hoặc hiểu $f$ như lớp tương đương hầu khắp nơi.
%ID BT_6_1_1_2_05
Tích thành tổng cho tích phân sine bằng
\[\frac{\sin((n-m)\pi/2)}{n-m}-\frac{\sin((n+m)\pi/2)}{n+m},\]
và cosine thay dấu trừ thành cộng. Vì $n-m$ chẵn nên $n+m$ cũng chẵn, hai tử bằng 0. Tuy nhiên đề cho $m,n\in\mathbb Z$ và chỉ loại $n=m$, chưa loại $n=-m$. Nếu $n=-m\ne0$, hai hàm cùng tần số và tích phân sine bằng $-\pi/2$, cosine bằng $\pi/2$, không bằng 0. Kết luận đúng với $n\ne\pm m$ (hoặc $m,n$ không âm khác nhau). Trường hợp chỉ số 0 xử lý trực tiếp; sine 0 đồng nhất 0.
%ID BT_6_1_1_2_06
Dùng $f=\sum b_n\sin k_nx$ và $f=a_0+\sum a_n\cos k_nx$, $k_n=n\pi/L$. Hệ số lần lượt $b_n=2L^{-1}\int_0^Lf\sin$, $a_n=2L^{-1}\int_0^Lf\cos$, $a_0=L^{-1}\int_0^Lf$.
(a) $b_n=2(1-(-1)^n)/(n\pi)$; cosine chỉ có $a_0=1$.
(b) $a_0=(1-\cos\alpha L)/(\alpha L)$,
\[b_n=\frac{2(-1)^nk_n\sin\alpha L}{L(\alpha^2-k_n^2)},\quad a_n=\frac{2\alpha(1-(-1)^n\cos\alpha L)}{L(\alpha^2-k_n^2)}.\]
Ở cộng hưởng lấy giới hạn; ví dụ $\alpha=k_N$ cho $b_N=1$, $b_n=0$ với $n\ne N$.
(c) $a_0=(e^{\alpha L}-1)/(\alpha L)$,
\[b_n=\frac{2k_n(1-(-1)^ne^{\alpha L})}{L(\alpha^2+k_n^2)},\quad a_n=\frac{2\alpha((-1)^ne^{\alpha L}-1)}{L(\alpha^2+k_n^2)}.\]
(d) $L=\pi$: $b_n=4(1-(-1)^n)/(\pi n^3)$, $a_0=\pi^2/6$, $a_n=-2(1+(-1)^n)/n^2$.
(e) $L=1$: $b_n=2[1+(-1)^n-2\cos(n\pi/2)]/(n\pi)$, $a_0=0$, $a_n=4\sin(n\pi/2)/(n\pi)$.

Sine là Fourier của thác triển lẻ chu kỳ $2L$, cosine là của thác triển chẵn. Tại điểm nội liên tục, cả hai về $f$; tại bước nhảy về trung bình hai phía. Sine luôn về 0 ở $0,L$, nên nếu $f$ không bằng 0 tại đó không có hội tụ đến giá trị biên. Cosine về giá trị biên một phía. Các cosine của (a)--(d) và sine của (b),(d) khi dữ liệu vanishes ở hai đầu có hội tụ đều nếu thác triển liên tục và các hệ số suy đủ nhanh; trường hợp bước nhảy (e) không đều. Các công thức hệ số chính xác ở trên cũng cho kiểm tra từng điều kiện biên.
